\uuid{36j0}
\titre{Limite}
\auteur{bodin ; hanani ; mzari}
\organisation{exo7}
\licence{CC BY-NC-SA 4.0}
\datecreate{2020-01-20}
\niveau{L1}
\module{Analyse}
\chapitre{Suite}
\sousChapitre{Autre}
\difficulte{1}
\theme{Suites, limite}
\exo7qcmid{lille-1.suites.005}

\contenu{
\begin{question}
\qclassification{121.00}
\qtags{motcle=limite}
Soit $\displaystyle u_n=3^n-2^n$ et $\displaystyle v_n=3^n-(-3)^n$. Quelles sont les bonnes réponses ?
\begin{answers}
    \bad{$\displaystyle \lim _{n\to +\infty}u_n=+\infty$ et $\displaystyle \lim _{n\to +\infty}v_n=+\infty$ }
    \bad{$\displaystyle \lim _{n\to +\infty}u_n=0$ et $\displaystyle \lim _{n\to +\infty}v_n=0$}
    \bad{$\displaystyle \lim _{n\to +\infty}u_n=0$ et $\displaystyle \lim _{n\to +\infty}v_n=-\infty$}
    \good{$\displaystyle \lim _{n\to +\infty}u_n=+\infty$ et $\displaystyle \lim _{n\to +\infty}v_n$ n'existe pas}
\end{answers}
\begin{explanations}
D'abord, $\displaystyle u_n=3^n\times \left[1-\left(\frac{2}{3}\right)^n\right]$. Or $\displaystyle \left(\frac{2}{3}\right)^n$ est le terme général d'une suite géométrique de limite $0$, donc $\displaystyle \lim _{n\to +\infty}u_n=+\infty\times (1-0)=+\infty$. Par ailleurs, $\displaystyle v_{2n}=0$ et $v_{2n+1}=2\times 3^{2n+1}$. Donc $\displaystyle \lim _{n\to +\infty}v_{2n}=0$ et $\displaystyle \lim _{n\to +\infty}v_{2n+1}=+\infty$. Le théorème des suites extraites implique que $(v_n)$ n'a pas de limite.
\end{explanations}
\end{question}
}
